% Felipe Muñoz Casasola
% This is free and unencumbered software released into the public domain.

% Anyone is free to copy, modify, publish, use, compile, sell, or
% distribute this software, either in source code form or as a compiled
% binary, for any purpose, commercial or non-commercial, and by any
% means.

% In jurisdictions that recognize copyright laws, the author or authors
% of this software dedicate any and all copyright interest in the
% software to the public domain. We make this dedication for the benefit
% of the public at large and to the detriment of our heirs and
% successors. We intend this dedication to be an overt act of
% relinquishment in perpetuity of all present and future rights to this
% software under copyright law.

% THE SOFTWARE IS PROVIDED "AS IS", WITHOUT WARRANTY OF ANY KIND,
% EXPRESS OR IMPLIED, INCLUDING BUT NOT LIMITED TO THE WARRANTIES OF
% MERCHANTABILITY, FITNESS FOR A PARTICULAR PURPOSE AND NONINFRINGEMENT.
% IN NO EVENT SHALL THE AUTHORS BE LIABLE FOR ANY CLAIM, DAMAGES OR
% OTHER LIABILITY, WHETHER IN AN ACTION OF CONTRACT, TORT OR OTHERWISE,
% ARISING FROM, OUT OF OR IN CONNECTION WITH THE SOFTWARE OR THE USE OR
% OTHER DEALINGS IN THE SOFTWARE.

% For more information, please refer to <https://unlicense.org/>

\documentclass{article}

\usepackage[spanish]{babel}
\usepackage[margin=3cm]{geometry}
\usepackage{microtype}
\usepackage[pdfusetitle]{hyperref}
\usepackage{minted}

\title{Copias de seguridad}
\author{Felipe Muñoz Casasola}
\date{22 de febrero de 2026}

\begin{document}
    \maketitle
    \clearpage

    \section{Instalación}
    Se instala como un contenedor a través de \texttt{podman}. Lo que se
    hace es fijar un \texttt{commit}, que se pueden consultar en los
    \textit{packages} de su repositorio oficial, para que no se esté
    actualizando continuamente y eso pueda provocar que se me rompa la
    configuración.

    \section{Actualización}
    Lo que se debe hacer es coger de su repositorio el archivo de configuración
    por defecto, copiarlo y cambiar una vez allí las opciones para que se
    adapten a las que yo quiero. De esta forma evito tener que guardar una
    \textit{cookie} para la configuración o modificar el enlace de la
    búsqueda. Una vez hecho esto, cambiar el \textit{commit} al que estoy
    anclado simplemente modificando el archivo en el que se inicia
    el contenedor.

    \section{Contenedor y autoarranque}
    Para hacer que se inicie desde el principio, se pueden seguir las instrucciones que haya actualizadas
    en la web y la página 5 del manual de la orden podman-systemd.unit. Ahora mismo lo que se deben
    utilizar son los \textit{quadlet}. Actualmente lo he generado con \texttt{podlet} que está todavía en beta, pero me
    servirá el ejemplo en caso de que hiciera falta actualizarlo.

    Para depurar los \textit{quadlet}, según la página 5 del manual anteriormente citado (
    \texttt{man 5 podman-systemd.unit}), puedo ejecutar: \texttt{/usr/lib/systemd/system-generators/podman-system-generator --dry-run}
    
    Dicho \textit{quadlet} es \texttt{$\sim$/.config/containters/systemd/searxng.container},
    y actualmente tiene el siguiente contenido:

    \inputminted[breaklines, breakanywhere, breaksymbolleft="", breaksymbolright="", breakanywheresymbolpre="", breakanywheresymbolpost=""]
    {systemd}{inc/searxng.container}

    \section{Configuración}
    Actualmente la configuración que tengo en SearxNG es la siguiente (es la
    que está en \texttt{$\sim$/.config/searxng/settings.yml}, que se redirige
    a \texttt{/etc/searxng} en el contenedor porque allí es donde realmente
    bsuca SearxNG sus archivos de configuración):

    \inputminted[breaklines, breakanywhere, breaksymbolleft="", breaksymbolright="", breakanywheresymbolpre="", breakanywheresymbolpost=""]
    {yaml}{inc/settings.yml}

    \section{Solucionar errores}
    Tengo que prestar atención a los errores haciendo un \texttt{status} con
    \texttt{systemctl --user}, para corregir los posibles errores en el archivo
    de configuración, si los hubiera. Nótese que a veces puede ocurrir que
    dé algún error diciendo que se está intentando asignar un \texttt{string} a
    un número o algo así. En dichos casos, me habré equivocado en algún sitio
    de la configuración y se me habrá escapado una letra en algún apartado.
\end{document}
